\section{Primitive roots and Euclidean division}

\subsection{Grosswald's conjecture and Erd\H{o}s problem 985}
For a prime $p$, write $g(p)$ for its least positive primitive root and $\widehat g(p)$ for its least prime primitive root. McGown, Trevi\~no and Trudgian establish the following implication in Theorem~1 of their 2016 paper:
\[
 \mathrm{GRH}_{\rm Dirichlet}\ \Longrightarrow\quad
 \begin{cases}
 g(p)<\sqrt{p}-2 & (p>409),\\
 \widehat g(p)<\sqrt{p}-2 & (p>2791).
 \end{cases}
\]
Their paper uses smoothed prime sums to detect a generator of the cyclic group $\mathbb F_p^\times$. For $p\nmid n$ the indicator of a primitive root is
\[
 \frac{\varphi(p-1)}{p-1}
 \sum_{d\mid p-1}\frac{\mu(d)}{\varphi(d)}
                 \sum_{\operatorname{ord}(\chi)=d}\chi(n).
\]
Here $\chi$ runs through Dirichlet characters modulo $p$. This formula follows by writing $n=u^k$ for a generator $u$ and applying inclusion--exclusion to $\gcd(k,p-1)=1$. Lemma~1 of the paper bounds the sums of $\Lambda(n)(1-n/x)$ and of $\Lambda(n)\chi(n)(1-n/x)$ for $n<x$. The first uses ordinary RH and the second uses Dirichlet GRH. A prime power generating $\mathbb F_p^\times$ has a prime base that also generates it. The authors' sieve and finite computations give the two cutoffs. Thus the repository's theorem supplies the analytic hypothesis of their theorem, including for the stronger prime-generator bound. The quantifier over characters here is the reason this use of GRH works. See \href{https://doi.org/10.7169/facm/2016.55.2.5}{McGown--Trevi\~no--Trudgian, Theorem~1}.

\paragraph{Corollary.}
For every odd prime $p$ there is a prime $q<p$ that generates $\mathbb F_p^\times$. This is the odd-prime form of \href{https://www.erdosproblems.com/985}{Erd\H{o}s problem~985}. For $p>2791$ take $q=\widehat g(p)$ and use the second inequality. For the remaining 405 odd primes, trial division enumerates $p$ and the distinct prime divisors $\ell$ of $p-1$. Test successive primes $q<p$ until
\[
      q^{(p-1)/\ell}\not\equiv1\pmod p
      \qquad\text{for every prime }\ell\mid p-1.
\]
The order of $q$ divides $p-1$ and any proper divisor divides $(p-1)/\ell$ for some such $\ell$, proving the test. The accompanying \texttt{checks/primitive\_roots.py} runs this test and independently enumerates all powers of each witness. Every odd prime in the range passes. The largest least witness is 53, attained only at 2791. The literal all-primes formulation has the elementary exception $p=2$. The strict endpoints of the stronger estimates also matter: $g(409)=21$ and $\widehat g(2791)=53$, so the respective estimates fail at their endpoints.

These deductions combine the repository with published implications and a small finite computation. For context, a 2026 large-prime result established the Grosswald bound for $p>1.77\times10^{54}$ without GRH. See \href{https://arxiv.org/abs/2608.20431}{Bagger, \emph{A note on least primitive roots}}.

\subsection{Which cyclic fields admit norm-Euclidean division?}
The ring of integers of $K$ is norm-Euclidean if, given $a,b\in\mathcal O_K$ with $b\ne0$, some $c\in\mathcal O_K$ satisfies $|N_{K/\mathbb Q}(a-bc)|<|N_{K/\mathbb Q}(b)|$. Lezowski and McGown classified these fields in two prime degrees assuming GRH. Their classification yields the following table upon supplying the repository's theorem.
\begin{center}
\begin{tabular}{c|l}
Degree of cyclic $K/\mathbb Q$ & Discriminants of norm-Euclidean $K$\\
\hline
$5$ & $11^4,\ 31^4,\ 41^4$\\
$7$ & $29^6,\ 43^6$\\
$19,31,37,43,47,59,67,71,73,79,97$ & none
\end{tabular}
\end{center}

\paragraph{Why the available GRH suffices.}
For a cyclic field $K/\mathbb Q$ of prime degree $\ell$ and conductor $f$, let $\chi$ be a generator of its group of Dirichlet characters. The abelian factorization is
\[
     \zeta_K(s)=\zeta(s)\prod_{j=1}^{\ell-1} L(s,(\chi^j)^*),
\]
where the star takes the associated primitive character. This convention includes the correct Euler factors at ramified primes. In the open critical strip the factors are holomorphic, so zero-freeness away from the critical line passes to their product. McGown explicitly specifies Dirichlet GRH as the default hypothesis of his conductor estimates, with individual $L(s,\chi)$ or $\zeta_K(s)$ indicated when needed. All these functions fall under the displayed factorization or the repository theorem directly. See \href{https://arxiv.org/abs/1102.2043}{McGown, \emph{Norm-Euclidean Galois fields and the Generalized Riemann Hypothesis}}, introduction and Theorems~3.1--3.2.

\paragraph{Where the classification is proved.}
Proposition~5.2 of Lezowski--McGown uses GRH bounds for character nonresidues and an auxiliary integer to bound the conductor of a norm-Euclidean field. Proposition~2.1 and Table~2.2 give $f<6\times10^{10}$ in degree five and $f<4\times10^{10}$ in degree seven. Proposition~2.2 supplies the finite exclusions, and Proposition~2.3 supplies the Euclidean minima of the surviving fields in these two degrees. Their computations and algebraic arguments give Theorems~1.3 and~1.4. The authors already proved the entire conditional implication; the contribution here is identifying precisely why the formal result meets its input. The extensive field computations are those of the published source. See \href{https://doi.org/10.1090/mcom/3169}{Lezowski--McGown, \emph{The Euclidean algorithm in quintic and septic cyclic fields}}.

Degree eleven illustrates the remaining arithmetic problem. Their Table~1.1 leaves conductors $23,67,331$ as candidates. Theorem~1.5 makes that list exhaustive under the same GRH. It supplies no positive norm-Euclidean certificate for an individual candidate. In degree three the complete classification already has a separate unconditional argument. See \href{https://arxiv.org/abs/2507.05862}{Bagger, Booker, Kerr, McGown, Starichkova and Trudgian}.
